\documentclass[11pt]{article}
\usepackage[margin=1in]{geometry}
\usepackage{amsmath,amssymb,amsthm}
\usepackage{graphicx}
\usepackage[hidelinks]{hyperref}

\newtheorem{theorem}{Theorem}
\newtheorem{lemma}{Lemma}
\newtheorem{corollary}{Corollary}

\newcommand{\Lip}{\operatorname{Lip}}
\newcommand{\LipF}{\operatorname{Lip}_{0\mathcal F}}
\newcommand{\ap}{\mathrm{a.p.}}
\newcommand{\pap}{p\text{-}\mathrm{a.p.}}

\title{The Lipschitz Approximation Properties in arXiv:2512.06317 Are Automatic}
\author{Automated discovery packet}
\date{June 8, 2026}

\begin{document}
\maketitle

\section*{Source question}

The final open-question list of arXiv:2512.06317 asks for examples of Banach
spaces without the Lipschitz approximation property or without the Lipschitz
$p$-approximation property, and asks what can be concluded about the reverse
implications in the displayed implication diagram.

\begin{center}
\includegraphics[width=\textwidth]{figures/open_problem_crop.png}
\end{center}

\section*{Result}

The question has a simple obstruction: with the definitions used in the
paper, the Lipschitz approximation property has no content unless one also
imposes a uniform bound on the Lipschitz constants of the approximating maps.

\begin{theorem}
Let $(M,0)$ be a pointed metric space, let $X$ be a Banach space, and let
$f\in \Lip_0(M,X)$. For every compact set $C\subset M$ and every
$\varepsilon>0$, there is a map $g\in\LipF(M,X)$ such that
\[
        \sup_{x\in C}\|g(x)-f(x)\|<\varepsilon .
\]
Consequently, every Banach space has the Lipschitz approximation property.
Since every relatively $p$-compact set is compact, every Banach space also has
the Lipschitz $p$-approximation property, for every $1\leq p<\infty$.
\end{theorem}

\begin{proof}
We use a finite-dimensional Lipschitz approximation of the identity on compact
subsets of the range. Let $K\subset X$ be compact and contain $0$. We first
show that, for every $\varepsilon>0$, there are a finite-dimensional subspace
$E\subset X$ and a Lipschitz map $R:K\to E$ such that $R(0)=0$ and
\[
        \sup_{y\in K}\|R(y)-y\|<\varepsilon .
\]
Put $\delta=\varepsilon/16$ and $A=\{y\in K:\|y\|\geq \delta\}$. If $A$ is
empty, take $R=0$. Otherwise choose a finite $\delta/2$-net
$a_1,\ldots,a_N$ in $A$ and set $E=\operatorname{span}\{a_1,\ldots,a_N\}$.
For $y\in A$ define
\[
 \theta_i(y)=\bigl(\delta-\|y-a_i\|\bigr)_+,\qquad
 D(y)=\sum_{i=1}^N\theta_i(y),
\]
and
\[
 q(y)=D(y)^{-1}\sum_{i=1}^N\theta_i(y)a_i .
\]
The net condition gives $D(y)\geq \delta/2$ on $A$, so $q$ is Lipschitz on
$A$. Moreover, whenever $\theta_i(y)>0$ we have $\|a_i-y\|<\delta$, hence
\[
        \|q(y)-y\|\leq \delta\qquad (y\in A).
\]
Let
\[
 \lambda(y)=\min\{1,\max\{0,(\|y\|-\delta)/\delta\}\}
\]
and define $R(y)=0$ if $\|y\|\leq\delta$, while $R(y)=\lambda(y)q(y)$ on
$A$. This is Lipschitz on $K$. Indeed, on $A$ it is a product and quotient of
Lipschitz functions with denominator bounded below. If $x\in A$ and
$\|y\|\leq\delta$, then
\[
 \|R(x)-R(y)\|=\lambda(x)\|q(x)\|
 \leq \frac{(\|x\|-\delta)_+}{\delta}(\|x\|+\delta)
 \leq \frac{M+\delta}{\delta}\|x-y\|,
\]
where $M=\sup_{z\in K}\|z\|$; the last inequality follows from
$\|x-y\|\geq\|x\|-\|y\|\geq \|x\|-\delta$. Thus the two Lipschitz pieces glue
with a finite constant. The approximation estimate is immediate. If
$\|y\|\leq\delta$, then $\|R(y)-y\|\leq\delta$.
If $\delta<\|y\|<2\delta$, then
\[
 \|R(y)-y\|\leq \|q(y)\|+\|y\|\leq (\|y\|+\delta)+\|y\|<5\delta.
\]
If $\|y\|\geq 2\delta$, then $R(y)=q(y)$ and $\|R(y)-y\|\leq\delta$. In all
cases the error is below $\varepsilon$.

Because $E$ is finite-dimensional, the Lipschitz map $R:K\to E$ extends to a
Lipschitz map $\widetilde R:X\to E$ with $\widetilde R(0)=0$: choose a basis
of $E$, extend the finitely many scalar coordinate functions by the McShane
extension theorem, and recombine them.

Now apply this construction to the compact set
$K=f(C)\cup\{0\}\subset X$. The map
\[
        g=\widetilde R\circ f
\]
is Lipschitz, vanishes at $0$, and has range contained in the finite-dimensional
space $E$, so $g\in\LipF(M,X)$. For $x\in C$,
\[
        \|g(x)-f(x)\|=\|\widetilde R(f(x))-f(x)\|<\varepsilon .
\]
Taking $M=X$ and $f=\operatorname{Id}_X$ proves the Lipschitz approximation
property for every Banach space. The $p$-version follows because the topology
$\tau_p$ tests approximation only on relatively $p$-compact sets, and these
sets are compact.
\end{proof}

\begin{corollary}
There are no Banach spaces without the Lipschitz approximation property, and
there are no Banach spaces without the Lipschitz $p$-approximation property,
with the unbounded definitions used in arXiv:2512.06317.
\end{corollary}

\begin{corollary}
The reverse implications
\[
        \Lip\text{-}\ap \Longrightarrow \ap
        \qquad\text{and}\qquad
        \Lip\text{-}p.\ap \Longrightarrow \ap
\]
are false. Indeed, let $E_0$ be Enflo's Banach space without the classical
approximation property. By the theorem, $E_0$ has both Lipschitz approximation
properties, but it does not have the classical approximation property.

For $1\leq p\leq2$, the reverse implication
\[
        \pap \Longrightarrow \ap
\]
is also false. In this range, the paper cites the Karn--Sinha theorem that
every Banach space has $\pap$. Hence Enflo's space again supplies a
counterexample.
\end{corollary}

\section*{Verification notes}

The proof is deterministic. The helper script
\texttt{code/check\_partition\_construction.py} mirrors the finite partition
construction on a compact sample in $\mathbb R^2$ and checks the advertised
error bound numerically. The script \texttt{code/make\_open\_problem\_crop.py}
recreates the source screenshot from \texttt{source\_paper.pdf}.

\begin{thebibliography}{9}
\bibitem{Mandal}
A. Mandal,
\emph{Approximation property in terms of Lipschitz maps via tensor product approach},
arXiv:2512.06317, 2025.

\bibitem{Enflo}
P. Enflo,
\emph{A counterexample to the approximation problem in Banach spaces},
Acta Mathematica 130 (1973), 309--317.

\bibitem{KarnSinha}
D. P. Sinha and A. K. Karn,
\emph{Compact operators whose adjoints factor through subspaces of $\ell_p$},
Studia Mathematica 150 (2002), 17--33.
\end{thebibliography}

\end{document}
